\subsection{Achieving Active Status} \label{VIII.1}

To achieve active status for the current
semester, a Michigan Gamma Tau Beta Pi member must satisfy one of the following:
\begin{enumerate}
    \item Be an officer, Distinguished Active, or Prestigious Active,
    \item Have initiated into the society in the current semester, or,
    \item Achieve a sufficient level of activity in the society, as defined by the following requirements:
    \begin{enumerate}
        \item Attend, at a minimum, two meetings, both of which must be voting meetings
(2nd Actives, 3rd Actives or Officer Elections),
        \item Complete two hours of TBP-sponsored service, and
        \item Attend at least one social event or serve as a mentor in the chapter mentorship
program.
    \end{enumerate}
\end{enumerate}
Members who have completed all requirements for active status apart from attendance
at voting meetings may be considered a Temporary Active if sufficient voting meetings
remain for the member to achieve active status for that semester. Temporary Active status
is only maintained through the duration of voting meeting for which it is granted.
In addition, any alumni member must request active status each semester, in accordance
with \href{http://www.tbp.org/off/ConstBylaw.pdf}{C-VI of the Constitution of the Tau Beta Pi Association}.


\subsection{Maintaining Active Status} \label{VIII.2}

Active status is maintained through the
end of the semester following the semester in which it is achieved, provided that the member
has attended at least one chapter meeting that semester. Those who leave campus may
maintain active status upon their return provided they achieved active status the semester
prior to leaving and attend at least one chapter meeting during the semester when they
return.


\subsection{Tracking Status} \label{VIII.3}

Active status is evaluated by the Membership Officer.
If an inactive member satisfies the above requirement mid-semester, they must also be
made active immediately upon request.


\subsection{Benefits of Active Status} \label{VIII.4}

Active status may be designated on some
chapter documents, and may be used in determining eligibility for certain chapter benefits
at the discretion of the officer corps. However, inactive members do not lose privileges
(other than voting) relating to standard meetings and socials.


\section{Enactment and Amendment} \label{IX}

\subsection{Enactment} \label{IX.1}

These Bylaws became effective immediately upon ratification
and supersede any other Bylaws of this chapter, existing at the time of ratification.


\subsection{Amendment} \label{IX.2}

The Chapter Bylaws may be amended upon the affirmative
vote of $\sfrac{2}{3}$ of the total active membership and an affirmative vote of a majority of the
Advisory Board. Such votes must occur within 14 days of each other.


\subsection{Appendix Amendment} \label{IX.3}

Appendix \ref{A} (Electee Requirements), Appendix
\ref{B} (Graduate Electee Requirements), Appendix \ref{C} (Alumni Electee Requirements), Appendix
\ref{D} (Distinguished Active Status Guidelines), and Appendix \ref{E} (Prestigious Active Status
Guidelines) may be modified at the discretion of the Officer Corps. Appendices \ref{F.4}, \ref{G.3},
and \ref{H} may be modified at the discretion of the Officer Corps per the procedures outlined
in Bylaws \hyperref[I.4.a.iii]{I.4.a.iii}, \hyperref[I.5.c]{I.5.c}, and \hyperref[I.4.a.iv]{I.4.a.iv} respectively.